% Updated 30 August 2026
\documentclass[10pt]{article}
\usepackage[margin=0.62in]{geometry}
\usepackage[T1]{fontenc}
\usepackage{lmodern,microtype}
\usepackage{amsmath,amssymb,amsthm,mathtools}
\usepackage{enumitem,booktabs,tabularx,array,multicol,xcolor}
\usepackage[bookmarks=false]{hyperref}
\hypersetup{colorlinks=true,linkcolor=blue!45!black,urlcolor=blue!45!black}
\setlength{\parindent}{0pt}
\setlength{\parskip}{3pt}
\setlist{leftmargin=*,itemsep=1pt,topsep=2pt}
\newcommand{\E}{\mathbb E}
\newcommand{\Prb}{\mathbb P}
\newcommand{\Var}{\operatorname{Var}}
\newcommand{\Cov}{\operatorname{Cov}}
\newcommand{\1}{\mathbf 1}
\newcommand{\cF}{\mathcal F}
\newcommand{\cG}{\mathcal G}
\newtheorem*{theorem}{Theorem}
\newtheorem*{lemma}{Lemma}
\newenvironment{keybox}{\begin{quote}\small\itshape}{\end{quote}}

\begin{document}
\begin{center}
{\LARGE\bfseries The Probabilistic Method\par}
\vspace{1mm}
{Hadi Zamani \qquad Examiner: Martin Tancer\par}
{\small 30 August 2026}
\end{center}

\section*{Assigned chapters}
\begin{itemize}
\item Chapter 6: Correlation Inequalities
\item Chapter 7: Martingales and Tight Concentration
\item Chapter 15: Codes, Games, and Entropy
\item Chapter 17: Graph Property Testing
\end{itemize}

\subsection*{Chapter map: main theorems and techniques}
\begingroup
\scriptsize
\setlength{\tabcolsep}{3.5pt}
\renewcommand{\arraystretch}{1.08}
\begin{center}
\begin{tabularx}{\textwidth}{@{}>{\raggedright\arraybackslash}p{0.08\textwidth}>{\raggedright\arraybackslash}X>{\raggedright\arraybackslash}X@{}}
\toprule
\textbf{Chapter} & \textbf{Main theorems and results} &
\textbf{Characteristic techniques}\\
\midrule
6 & Ahlswede--Daykin Four Functions; FKG and Harris--Kleitman correlation;
the XYZ theorem for random linear extensions. & Finite distributive lattices;
log-supermodularity; increasing/decreasing events; complement one event for negative
correlation; limiting perturbations of random labels.\\
7 & Azuma--Hoeffding/McDiarmid; Shamir--Spencer concentration of chromatic
number; Talagrand convex-distance inequality; Kim--Vu polynomial concentration.
& Doob martingales and carefully chosen exposures; exponential moments;
Lipschitz bounds plus small certificates; formal derivatives and expectation hierarchy.\\
15 & Binary symmetric channel capacity $1-H(p)$; Hamming-sphere/ChipLiar bounds;
balancing-game value $\Theta(\sqrt{n\log n})$; Shearer's entropy inequality and its
projection corollary. & Random coding and entropy estimates for Hamming balls;
random adversarial strategies; potential functions; entropy chain rule, conditioning,
and submodularity.\\
17 & One-sided testing of $k$-colourability by an
$O(k\log k/\varepsilon^2)$-vertex sample (the stated $k=3$ bound is
$40/\varepsilon^2$); triangle-removal testing; every hereditary dense-graph property
is one-sided testable. & Random induced samples; colour-list restriction trees;
Szemer\'edi regularity and removal; induced-removal for forbidden induced subgraphs.\\
\bottomrule
\end{tabularx}
\end{center}
\endgroup

\section*{Theorem inventory}
\begingroup
\small

\subsection*{Short proof-hint index}
Use these as the first sentence of an oral proof; the detailed arguments appear
in the corresponding chapter sections below.
\begin{itemize}
\item \textbf{Chapter 6:} verify the pointwise join/meet inequality, sum it by
the Four Functions Theorem, and specialise the log-supermodular measure to
increasing events (FKG/Harris--Kleitman).
\item \textbf{Chapter 7:} reveal independent coordinates to form a Doob
martingale; use bounded differences for Azuma, and combine Lipschitz control
with a small certificate for Talagrand's convex-distance bound.
\item \textbf{Chapter 15:} estimate Hamming-ball volume by binary entropy,
then use random coding/union bounds; for liar games, track the sphere-bound
potential after every answer.
\item \textbf{Chapter 17:} sample an induced subgraph; if the graph is far from
the property, convert that distance into many small witnesses, using removal
or induced-removal when the forbidden family is not fixed.
\end{itemize}

\subsection*{Chapter 6: correlation inequalities}
\begin{itemize}
\item \textbf{Theorem 6.1.1 (Four Functions).} If nonnegative functions on
$2^N$ satisfy
$\alpha(A)\beta(B)\leq\gamma(A\cup B)\delta(A\cap B)$ pointwise, then the
corresponding inequality holds after summing over any two families.
\item \textbf{Corollary 6.1.2.} The Four Functions Theorem holds on every
finite distributive lattice, with union/intersection replaced by join/meet.
\item \textbf{Corollary 6.1.3.} For $X,Y$ in a finite distributive lattice,
$|X||Y|\leq|X\vee Y||X\wedge Y|$.
\item \textbf{Corollary 6.1.4.} Every set family $\mathcal A$ satisfies
$|\mathcal A-\mathcal A|\geq|\mathcal A|$.
\item \textbf{Theorem 6.2.1 (FKG).} Under a log-supermodular weight on a
finite distributive lattice, two nonnegative increasing functions are positively
correlated.
\item \textbf{Proposition 6.3.1 (Harris--Kleitman, uniform case).} Two
increasing events, or two decreasing events, are positively correlated; an
increasing and a decreasing event are negatively correlated.
\item \textbf{Theorem 6.3.2 (Harris--Kleitman, product measure).} The same
three correlation inequalities hold for arbitrary independent coordinate
probabilities $p_1,\ldots,p_n$.
\item \textbf{Theorem 6.3.3.} Increasing/decreasing random-graph properties
obey the corresponding positive/negative correlation inequalities in $G(n,p)$.
\item \textbf{Theorem 6.4.1 (XYZ theorem).} In a uniformly random linear
extension of a finite poset,
$\Prb[a_1<a_2\text{ and }a_1<a_3]\geq
\Prb[a_1<a_2]\Prb[a_1<a_3]$.
\end{itemize}

\subsection*{Chapter 7: martingales and tight concentration}
\begin{itemize}
\item \textbf{Theorem 7.2.1 (Azuma).} If $X_0=0,\ldots,X_m$ is a martingale
with $|X_{i+1}-X_i|\leq1$, then
$\Prb[X_m>\lambda\sqrt m]<e^{-\lambda^2/2}$.
\item \textbf{Corollary 7.2.2.} For $X_0=c$ under the same bound,
$\Prb[|X_m-c|>\lambda\sqrt m]<2e^{-\lambda^2/2}$.
\item \textbf{Theorem 7.2.3.} Edge-Lipschitz functions have unit differences
in the edge-exposure martingale; vertex-Lipschitz functions have unit
differences in the vertex-exposure martingale.
\item \textbf{Theorem 7.2.4 (Shamir--Spencer).} For $G\sim G(n,p)$,
$\chi(G)$ lies within $\lambda\sqrt{n-1}$ of its mean except with probability
at most $2e^{-\lambda^2/2}$.
\item \textbf{Lemma 7.3.1.} The auxiliary edge-disjoint clique-packing variable
$Y$ satisfies $\E Y\geq(1+o(1))n^2/(2k^4)$.
\item \textbf{Theorem 7.3.2.} For the stated $k\sim2\log_2 n$ regime,
$\Prb[\omega(G)<k]$ is at most
$\exp\{-\Omega(n^2/\ln^8 n)\}$.
\item \textbf{Theorem 7.3.3.} If $p=n^{-\alpha}$ with fixed
$\alpha>5/6$, then for some integer $u=u(n,p)$,
$u\leq\chi(G(n,p))\leq u+3$ asymptotically almost surely.
\item \textbf{Lemma 7.3.4.} Under the hypotheses of Theorem 7.3.3, every
$c\sqrt n$-vertex induced subgraph is $3$-colourable asymptotically almost
surely.
\item \textbf{Theorem 7.4.1.} A function Lipschitz relative to a gradation has
Doob-martingale differences bounded by $1$.
\item \textbf{Theorem 7.4.2.} Such a length-$m$ gradation gives both Azuma
tails $e^{-\lambda^2/2}$ at distance $\lambda\sqrt m$ from the mean.
\item \textbf{Theorem 7.4.3.} An adaptive exposure strategy whose every line
has total conditional variance at most $\sigma^2$ gives an essentially Gaussian
tail $2e^{-\alpha^2/(2(1+\varepsilon))}$ in its stated range.
\item \textbf{Theorem 7.5.1.} The number $L(g)$ of values missed by a uniformly
random map $g:[n]\to[n]$ is concentrated within $\lambda\sqrt n+1$ of $n/e$.
\item \textbf{Theorem 7.5.2.} For independent random signs and
$\|v_i\|\leq1$, the norm $\|\sum_i\varepsilon_i v_i\|$ has two-sided
subgaussian deviations of scale $\sqrt n$ about its mean.
\item \textbf{Theorem 7.5.3.} If $A\subseteq\{0,1\}^n$ has density at least
$\varepsilon$, then a Hamming neighbourhood of radius
$2\sqrt{2n\ln(1/\varepsilon)}$ covers density at least $1-\varepsilon$.
\item \textbf{Theorem 7.6.1.} Talagrand's convex distance from $x$ to $A$
equals the Euclidean distance from the origin to the convex hull of the
admissible difference vectors $U(A,x)$.
\item \textbf{Theorem 7.6.2 (Talagrand).} On a product space,
$\int\exp(\rho(A,x)^2/4)\,dx\leq1/\Prb[A]$.
\item \textbf{Theorem 7.7.1.} A $1$-Lipschitz, $f$-certifiable variable $X$
satisfies
$\Prb[X\leq b-t\sqrt{f(b)}]\Prb[X\geq b]\leq e^{-t^2/4}$.
\item \textbf{Theorem 7.8.1 (Kim--Vu).} A positive degree-$k$ polynomial of
independent Bernoulli variables concentrates around its mean at scale
$a_k(EE')^{1/2}\lambda^k$, with failure probability at most
$d_ke^{-\lambda}n^{k-1}$.
\end{itemize}

\subsection*{Chapter 15: codes, games, and entropy}
\begin{itemize}
\item \textbf{Theorem 15.1.1 (Shannon).} For a binary symmetric channel with
error probability $p<1/2$, rates arbitrarily close to $1-H(p)$ from below are
achievable with arbitrarily small decoding error.
\item \textbf{Theorem 15.1.2.} The same capacity lower bound is achievable by
binary linear (group) codes.
\item \textbf{Theorem 15.2.1.} In the ChipLiar game, if
$\sum_{i=0}^k x_i B(q,i)>1$, then Carole has a winning strategy.
\item \textbf{Corollary 15.2.2 (sphere bound).} If
$n>2^q/\sum_{i=0}^k\binom qi$, then Carole wins the $(n,q,k)$ liar game.
\item \textbf{Theorem 15.3.1.} If $\sum_i x_i2^{-i}<1$, Carole wins the
$(x_1,\ldots,x_k)$ tenure game.
\item \textbf{Lemma 15.3.2.} Chips of total weight at least $1$ can be split
into two parts, each of weight at least $1/2$.
\item \textbf{Theorem 15.3.3.} If $\sum_i x_i2^{-i}\geq1$, Paul wins the
tenure game; together with Theorem 15.3.1 this is an exact criterion.
\item \textbf{Theorem 15.4.1.} If
$\Prb[|S_n|>\alpha]<1/n$, then the balancing-vector game has
$\operatorname{VAL}(n)\leq\alpha$.
\item \textbf{Corollary 15.4.2.}
$\operatorname{VAL}(n)=O(\sqrt{n\ln n})$.
\item \textbf{Theorem 15.4.3.} If
$\Prb[|S_n|>\alpha]>cn^{-1/2}$ for the stated absolute constant $c$, then
$\operatorname{VAL}(n)>\alpha$.
\item \textbf{Corollary 15.4.4.}
$\operatorname{VAL}(n)=\Theta(\sqrt{n\ln n})$.
\item \textbf{Section 15.5.} There is no numbered theorem; the balancing game
identifies the optimal worst-case guarantee for nonadaptive online discrepancy.
\item \textbf{Theorem 15.6.1.} For every fixed $k\geq1$, the half-liar-game
capacity satisfies $A_k(q)\sim2^{q+k}/\binom qk$.
\item \textbf{Lemma 15.7.1.} Entropy is bounded by log-support size, is
monotone under adjoining variables, is subadditive, and decreases under
conditioning.
\item \textbf{Proposition 15.7.2.} Joint entropy is at most the sum of the
coordinate entropies: $H(X_1,\ldots,X_n)\leq\sum_iH(X_i)$.
\item \textbf{Corollary 15.7.3.} If $p_i$ is the fraction of members of a set
family containing $i$, then
$|\mathcal F|\leq2^{\sum_iH(p_i)}$.
\item \textbf{Proposition 15.7.4 (Shearer).} If each coordinate belongs to at
least $k$ members of $\mathcal G$, then
$kH(X)\leq\sum_{G\in\mathcal G}H(X_G)$.
\item \textbf{Corollary 15.7.5.} For a family of vectors and coordinate
projections $\mathcal F_i$, Shearer's condition gives
$|\mathcal F|^k\leq\prod_i|\mathcal F_i|$.
\item \textbf{Corollary 15.7.6 (Loomis--Whitney).} The volume of a body obeys
$\operatorname{Vol}(B)^{n-1}\leq\prod_i\operatorname{Vol}(B_i)$.
\item \textbf{Corollary 15.7.7.} The projection inequality of Corollary 15.7.5
specializes to families of subsets.
\item \textbf{Corollary 15.7.8.} A family of labelled graphs in which every two
members share a triangle has size less than
$2^{\binom t2}/4$.
\end{itemize}

\subsection*{Chapter 17: graph property testing}
\begin{itemize}
\item \textbf{Theorem 17.2.1.} If $0<\varepsilon<0.1$,
$n>400/\varepsilon^3$, and an $n$-vertex graph is
$\varepsilon$-far from $3$-colourable, then a uniformly random induced
subgraph on $40/\varepsilon^2$ vertices is $3$-colourable with probability at
most $0.1$.
\item \textbf{Claims 17.2.2--17.2.3.} The colour-list extension bound converts
distance from $3$-colourability into a set of at least $\varepsilon n/2$
vertices that force progress in the restriction tree.
\item \textbf{Claims 17.2.4--17.2.6.} The restriction tree has depth at most
$3/\varepsilon$; once all leaves are terminal the sample is not
$3$-colourable; after $40/\varepsilon^2$ samples this occurs with probability
at least $0.9$.
\item \textbf{Lemma 17.3.1 (triangle removal).} Every graph
$\varepsilon$-far from triangle-free contains at least
$\delta(\varepsilon)n^3$ triangles.
\item \textbf{Corollary 17.3.2.} Triangle-freeness has a one-sided
constant-query tester.
\item \textbf{Lemma 17.4.1.} A uniform induced-removal statement for a family
$\mathcal F$ implies one-sided testability of induced-$\mathcal F$-freeness.
\item \textbf{Theorem 17.4.2 (Alon--Shapira).} The required induced-removal
functions exist for every finite or infinite forbidden family $\mathcal F$.
\item \textbf{Theorem 17.4.3.} Every hereditary dense-graph property is
testable with one-sided error.
\item \textbf{Corollary 17.4.4.} If a graph is $\varepsilon$-far from a
hereditary property $\mathcal P$, it has a bounded-size induced subgraph, of
size at most $W_{\mathcal P}(\varepsilon)$, that violates $\mathcal P$.
\item \textbf{Theorem 17.4.5.} A graph property has an oblivious one-sided
tester if and only if it is semi-hereditary.
\end{itemize}

\endgroup

\section{Chapter 6: correlation inequalities}

\subsection{Lattices, monotonicity, and the Four Functions Theorem}

A finite lattice $(L,\leq)$ has a meet $x\wedge y$ and join $x\vee y$.
It is distributive when meet and join distribute over each other.  The main
example is $L=2^{[n]}$ ordered by inclusion, with
$A\wedge B=A\cap B$ and $A\vee B=A\cup B$.
A function $f:L\to\mathbb R$ is increasing if $x\leq y$ implies
$f(x)\leq f(y)$.

For nonnegative $\alpha,\beta,\gamma,\delta:L\to\mathbb R_{\geq0}$, suppose
\[
 \alpha(x)\beta(y)\leq\gamma(x\vee y)\delta(x\wedge y)
 \qquad\text{for all }x,y\in L.                       \tag{4F-local}
\]
Writing $\alpha(\mathcal A)=\sum_{x\in\mathcal A}\alpha(x)$ and
\[
 \mathcal A\vee\mathcal B=\{x\vee y:x\in\mathcal A,y\in\mathcal B\},
 \qquad
 \mathcal A\wedge\mathcal B=\{x\wedge y:x\in\mathcal A,y\in\mathcal B\},
\]
the Ahlswede--Daykin Four Functions Theorem gives
\[
 \alpha(\mathcal A)\beta(\mathcal B)
 \leq\gamma(\mathcal A\vee\mathcal B)
       \delta(\mathcal A\wedge\mathcal B).             \tag{4F-global}
\]
The striking point is that a pointwise lattice inequality survives summation.
The proof is an induction on the lattice/ground-set size, splitting according
to one coordinate and applying the hypothesis to the induced functions.

\subsection{FKG and Harris--Kleitman}

A nonnegative weight $\mu:L\to\mathbb R_{\geq0}$ is
\emph{log-supermodular} if
\[
 \mu(x\vee y)\mu(x\wedge y)\geq\mu(x)\mu(y).
\]
After normalization, regard $\mu$ as a probability distribution.  The FKG
inequality states that for increasing $f,g:L\to\mathbb R$,
\[
 \E_\mu[fg]\geq\E_\mu[f]\,\E_\mu[g].                  \tag{FKG}
\]
For increasing events $A,B$, take indicator functions:
\[
 \Prb_\mu(A\cap B)\geq\Prb_\mu(A)\Prb_\mu(B).
\]
Proof route: combine log-supermodularity of $\mu$ and monotonicity of $f,g$ to
form four functions satisfying (4F-local), apply Four Functions, and normalize.

For independent coordinates, the product measure on $2^{[n]}$ has
\[
 \mu(S)=\prod_{i\in S}p_i\prod_{i\notin S}(1-p_i)
\quad\Longrightarrow\quad
\mu(A\cup B)\mu(A\cap B)=\mu(A)\mu(B).
\]
Therefore increasing events are positively correlated (Harris--Kleitman).
If $A$ is increasing and $D$ is decreasing, then $D^c$ is increasing and
\[
 \Prb(A\cap D)\leq\Prb(A)\Prb(D).
\]

\textbf{Graph-property application.}  In $G(n,p)$, containing a Hamilton cycle
and being connected are increasing; planarity and being triangle-free are
decreasing.  Thus two increasing properties are positively correlated, while
an increasing and a decreasing property are negatively correlated.
Independence is false in general even though the edges themselves are independent.

\subsection{The XYZ theorem for linear extensions}

Choose a linear extension of a finite poset $P$ uniformly.  For any elements
$x,y,z$ (where $x<y$ denotes that $x$ precedes $y$ in the random extension),
\[
 \Prb[x<y\ \wedge\ x<z]\geq
 \Prb[x<y]\Prb[x<z].                                  \tag{XYZ}
\]
The proof embeds order-preserving integer labelings into a carefully chosen
finite distributive lattice.  The set of labelings has a log-supermodular
indicator, while the two order events are increasing.  Apply FKG and let the
label range tend to infinity; ties become negligible, leaving uniform linear
extensions.  The lesson is not that arbitrary order events correlate positively:
the lattice order and monotonicity must be verified.

\begin{keybox}
Correlation checklist: identify the lattice; determine which events/functions
are increasing; verify log-supermodularity; apply FKG; reverse the sign by
complementing exactly one increasing event.  Never infer correlation merely from
a verbal intuition that two events are both likely.
\end{keybox}

\section{Chapter 7: martingales and tight concentration}

\subsection{Definitions and the Doob construction}

For a filtration
$\mathcal F_0\subseteq\cdots\subseteq\mathcal F_m$, a process $(X_i)$ is a
martingale if $X_i$ is $\mathcal F_i$-measurable, integrable, and
\[
 \E[X_{i+1}\mid\mathcal F_i]=X_i.
\]
Given any integrable random variable $Y$, exposing random choices one at a time
gives the Doob martingale
\[
 X_i=\E[Y\mid\mathcal F_i],\qquad X_0=\E Y,\quad X_m=Y.
\]
The method separates the work into two parts: choose an exposure revealing $Y$,
then bound how much each new piece of information changes the conditional
expectation.

\subsection{Azuma--Hoeffding and bounded differences}

If $D_i=X_i-X_{i-1}$ and $|D_i|\leq c_i$ almost surely, Hoeffding's lemma gives
\[
 \E[e^{\lambda D_i}\mid\mathcal F_{i-1}]
 \leq e^{\lambda^2c_i^2/2}.
\]
Iterating conditional expectation and applying Markov to
$e^{\lambda(X_m-X_0)}$, then optimizing
$\lambda=t/\sum_i c_i^2$, yields
\[
 \Prb[X_m-X_0\geq t]\leq
 \exp\!\left(-\frac{t^2}{2\sum_i c_i^2}\right),
\]
\[
 \Prb[|X_m-X_0|\geq t]\leq
 2\exp\!\left(-\frac{t^2}{2\sum_i c_i^2}\right).       \tag{Azuma}
\]

McDiarmid's bounded-differences inequality is the product-space corollary:
if $Y=f(Z_1,\ldots,Z_m)$ for independent coordinates and changing coordinate
$i$ changes $f$ by at most $c_i$, then the same two-sided tail bound holds
around $\E Y$.

\subsection{Two exposure choices in random graphs}

\textbf{Edge exposure.}  Reveal the $\binom n2$ potential edges.  Use this when
one edge changes the statistic by $O(1)$.  For instance, the maximum number of
edge-disjoint $K_r$'s is edge-Lipschitz, while the total number of $K_r$'s is
not: one edge may lie in many cliques.

\textbf{Vertex exposure.}  At step $i$, reveal all still-unrevealed edges from
vertex $i$.  If two graphs differ only in the edges incident with one vertex,
their chromatic numbers differ by at most one, because deleting that vertex
leaves the same graph.  Hence for $G\sim G(n,p)$ and
$c=\E[\chi(G)]$,
\[
 \Prb\!\left(|\chi(G)-c|>\lambda\sqrt{n-1}\right)
 <2e^{-\lambda^2/2}.                                  \tag{Shamir--Spencer}
\]
The coarser coordinate block produces a much stronger bound than treating
all $\binom n2$ edges separately.

\subsection{General gradations and certificate-sensitive bounds}

A gradation
$\varnothing=B_0\subset B_1\subset\cdots\subset B_m=B$
exposes a random map one block at a time.  If a statistic is Lipschitz relative
to the gradation, then the induced Doob differences are bounded and Azuma
applies.  In sparse product spaces, a query/certificate analysis can replace
the crude number of all coordinates by the variance proxy
\[
 \sigma^2=\sum_i p_i(1-p_i)c_i^2,
\]
provided an adaptive procedure determines the statistic with bounded query
influence.  State the hypotheses before using the improved exponent.

\subsection{Talagrand's convex-distance inequality}

Let $\Omega=\prod_{i=1}^n\Omega_i$ have product measure.  For $A\subseteq\Omega$
and $x\in\Omega$, define
\[
 d_T(x,A)=\sup_{\|\alpha\|_2=1}\ 
 \inf_{y\in A}\sum_{i:x_i\neq y_i}|\alpha_i|.
\]
Talagrand's inequality states
\[
 \Prb(A)\Prb\{x:d_T(x,A)\geq t\}\leq e^{-t^2/4}.        \tag{Talagrand}
\]
Ordinary Hamming distance counts changed coordinates equally; $d_T$ lets the
worst unit vector weight important coordinates.  An equivalent exponential
moment form is
\[
 \E\!\left[e^{d_T(X,A)^2/4}\right]\leq\frac1{\Prb(A)}.
\]

Application pattern: show $Y$ is 1-Lipschitz and that $Y\geq s$ can be certified
by at most $f(s)$ coordinates.  If $Y$ is far above a median, every nearby point
must preserve enough certificate coordinates, forcing large convex distance
from the lower-level set.  Talagrand can then yield a fluctuation scale
$O(\sqrt{f(\operatorname{med}Y)})$, much sharper than Azuma's
$O(\sqrt n)$ when certificates are small.

More explicitly, let $M$ be a median and $A=\{Y\leq M\}$.  If $Y(x)\geq M+t$
and a certificate for $x$ uses at most $f(Y(x))$ coordinates, then any $y\in A$
must disagree with $x$ on at least $t$ certificate coordinates (under the usual
unit-Lipschitz certificate hypothesis).  Give those certificate coordinates equal
weight $1/\sqrt{f(Y(x))}$.  This exhibits
\[
 d_T(x,A)\geq \frac{t}{\sqrt{f(Y(x))}}.
\]
Substitution into Talagrand's inequality converts certificate size directly into a
tail bound.  The proof step to remember is the construction of the weight vector;
quoting ``small certificates'' without it leaves the concentration argument incomplete.

\subsection{Kim--Vu polynomial concentration}

Let $H$ be a weighted hypergraph of rank at most $k$, let $t_i$ be independent
indicators, and
\[
 Y=\sum_{e\in E(H)}w_e\prod_{i\in e}t_i,\qquad \mu=\E Y.
\]
For $A\subseteq V(H)$ let $Y_A$ be the formal partial derivative with respect to
the variables indexed by $A$, put
\[
 E_j=\max_{|A|=j}\E Y_A,\qquad
 E'=\max_{1\leq j\leq k}E_j,\qquad E=\max\{\mu,E'\}.
\]
The Kim--Vu theorem gives, for $\lambda>1$,
\[
 \Prb\!\left(|Y-\mu|>a_k(EE')^{1/2}\lambda^k\right)
 <d_k e^{-\lambda}n^{k-1},                            \tag{Kim--Vu}
\]
where one may take $a_k=8^k\sqrt{k!}$ and $d_k=2e^2$.
For fixed $k$, choose $\lambda\gg\log n$ so the exponential beats the polynomial
factor.  Computing the derivative expectations $E_j$ is the combinatorial core;
they measure the largest expected extension count after conditioning on a
partial pattern.

\begin{keybox}
Concentration checklist: choose coordinates/exposure; prove the exact Lipschitz
or difference bound; locate the center (expectation or median); compare the
desired deviation with the resulting scale.  Azuma may fail because the
exposure is poor, not because the variable is genuinely unconcentrated.
\end{keybox}

\section{Chapter 15: codes, games, and entropy}

\subsection{Binary entropy and volumes}

For $p\in(0,1)$, the binary entropy is
\[
 H(p)=-p\log_2p-(1-p)\log_2(1-p).
\]
For fixed $0<p<1/2$,
\[
 \sum_{i\leq pn}\binom ni=2^{H(p)n+o(n)}.
\]
This is the exponential size of a Hamming ball of relative radius $p$.
Whenever a coding bound contains $1-H(p)$, it comes from the exponent
$2^n/2^{H(p)n}$: the ambient cube divided by the effective error-ball volume.

\subsection{Shannon's theorem by random coding}

In the binary symmetric channel $\operatorname{BSC}(p)$, every transmitted bit
is independently flipped with probability $p<1/2$.  A code with $M=2^{Rn}$
length-$n$ codewords has rate $R$.  Shannon's theorem says that for every
$\varepsilon>0$ there are schemes of rate
\[
 R>1-H(p)-\varepsilon
\]
with decoding error below $\varepsilon$ for sufficiently large $n$.

Proof skeleton: select codewords randomly and decode to the unique codeword
within Hamming distance $(p+\eta)n$ when one exists.  Chernoff bounds make the
actual error pattern atypically large with probability $o(1)$.  The ball bound
gives at most $2^{H(p+\eta)n}$ plausible received strings around a competitor.
A union bound over $M-1$ competitors is exponentially small when
\[
 R+H(p+\eta)<1.
\]
Averaging supplies a deterministic code.  Random linear/group codes use the
same collision calculation while giving a succinct algebraic encoder.
The converse direction says a rate above $1-H(p)+\varepsilon$ must retain
significant error; capacity is therefore $1-H(p)$.

\subsection{The liar game and the sphere bound}

Paul must identify one of $n$ possibilities using $q$ yes/no questions while
Carole may lie at most $k$ times.  A truthful answer string has
\[
 V(q,k)=\sum_{j=0}^k\binom qj
\]
strings within $k$ lies.  Distinguishability therefore requires the Hamming
sphere bound
\[
 nV(q,k)\leq2^q.                                      \tag{sphere bound}
\]

In the ChipLiar formulation, let $x_i$ be the number of chips that can tolerate
$i$ further lies and put
\[
 B(q,i)=2^{-q}\sum_{j=0}^i\binom qj.
\]
If
\[
 \sum_{i=0}^k x_iB(q,i)>1,
\]
then Carole wins.  Proof: fix any strategy of Paul and let Carole choose each
answer by a fair coin.  A chip starting at level $i$ survives with probability
$B(q,i)$, so the expected number of survivors exceeds one.  Some answer sequence
therefore leaves at least two candidates.  This is a clean example of using a
random strategy to prove that an adversarial player has a deterministic winning
strategy.

\subsection{Balancing, tenure, and online viewpoints}

In the balancing game, Pusher proposes a sign vector and Chooser selects its
global sign.  A potential based on random-walk tail probabilities makes the two
successor potentials average to the current one.  Chooser selects the smaller
potential for an upper bound; Pusher selects coordinates to keep the two
successors close for a lower bound.  This yields the scale
\[
 \operatorname{VAL}(n)=\Theta(\sqrt{n\log n}).
\]
The nonadaptive/online discrepancy interpretation is that columns arrive one at
a time and their signs must be fixed irrevocably.  The half-liar game changes
the allowed direction of a lie, so ordinary symmetric Hamming balls no longer
describe the state space; the weighted decision-tree argument must respect the
asymmetry.  Across these games the invariant is the same: define a potential
whose average over the opponent's two moves is controlled.

\subsection{Entropy identities and Shearer's inequality}

For a finite random variable $X$,
\[
 H(X)=-\sum_x\Prb(X=x)\log_2\Prb(X=x),\qquad
 H(X\mid Y)=H(X,Y)-H(Y).
\]
Core facts:
\[
 0\leq H(X)\leq\log_2|\operatorname{supp}X|,
\quad H(X,Y)=H(Y)+H(X\mid Y),
\]
\[
 H(X\mid Y,Z)\leq H(X\mid Y),
\qquad
 H(X_1,\ldots,X_n)\leq\sum_iH(X_i).
\]
Equality in subadditivity holds for mutually independent coordinates.

Let $X=(X_1,\ldots,X_n)$ and let $\mathcal A$ be a multiset of subsets of
$[n]$ such that every coordinate belongs to at least $r$ members.  Shearer's
inequality is
\[
 rH(X)\leq\sum_{A\in\mathcal A}H(X_A).                 \tag{Shearer}
\]
One proof repeatedly replaces two sets by their union and intersection, using
submodularity of entropy,
\[
 H(X_{A\cup B})+H(X_{A\cap B})
 \leq H(X_A)+H(X_B),
\]
until one copy of $[n]$ can be removed and induction on $r$ applies.

If $\mathcal F\subseteq S_1\times\cdots\times S_n$, choose $X$ uniformly from
$\mathcal F$.  For the projection $\mathcal F_A$ onto coordinates $A$,
$H(X)=\log_2|\mathcal F|$ and $H(X_A)\leq\log_2|\mathcal F_A|$.  Therefore
\[
 |\mathcal F|^r\leq\prod_{A\in\mathcal A}|\mathcal F_A|.
\]
This entropy-to-counting conversion is a standard oral-exam proof pattern.

\section{Chapter 17: graph property testing}
\begingroup
\small

\subsection{Dense-graph model and definitions}

The input is an $n$-vertex graph accessed through adjacency queries:
"is $uv$ an edge?"  A graph is $\varepsilon$-far from property $\mathcal P$
if at least $\varepsilon n^2$ edge additions/deletions are needed to put it in
$\mathcal P$.  A tester makes $q(\varepsilon)$ queries, independent of $n$:
\[
\begin{array}{ll}
G\in\mathcal P &\Longrightarrow \Prb[\text{accept}]\geq2/3,\\
G\text{ is }\varepsilon\text{-far from }\mathcal P
 &\Longrightarrow \Prb[\text{reject}]\geq2/3.
\end{array}
\]
One-sided error means every $G\in\mathcal P$ is accepted with probability one;
rejection exhibits an actual forbidden witness.  An oblivious tester samples a
constant number of vertices and decides from their induced subgraph.

\begin{keybox}
This is the dense adjacency-query model.  In the bounded-degree model, distance
is normalized by $\varepsilon dn$ edge edits and queries ask for neighbors.
Theorems and query complexities do not transfer automatically between the two
models.
\end{keybox}

\subsection{Testing $k$-colourability}

The one-sided tester samples a uniform vertex set $R$ and accepts iff $G[R]$ is
$k$-colourable.  If $G$ is $k$-colourable, every induced subgraph is, so
completeness is perfect.  The hard direction is a small-witness theorem:
if $G$ is $\varepsilon$-far from $k$-colourable, a random induced subgraph on
\[
 O\!\left(\frac{k\log k}{\varepsilon^2}\right)
\]
vertices is non-$k$-colourable with high constant probability.  For the
book's $k=3$ presentation, when $0<\varepsilon<0.1$ and
$n>400/\varepsilon^3$, a sample of $40/\varepsilon^2$ vertices is
3-colourable with probability at most $0.1$.

Proof architecture for $k=3$: maintain feasible colour lists relative to a
partially coloured sample.  Either many vertices are already colourless, so the
sample soon hits a direct obstruction, or many vertices are \emph{restricting}:
whatever feasible colour they receive, they remove that colour from at least
$\varepsilon n$ other lists.  Each restriction lowers the total list mass,
so a branch has depth at most $3/\varepsilon$.  A bounded branching tree plus a
sampling tail bound shows that all candidate colourings acquire an obstruction.

Here is the quantitative engine.  For a proper partial colouring
$\phi:S\to\{1,2,3\}$, let $L_\phi(v)$ be the colours not yet forbidden by
neighbours of $v$ in $S$, and let $U$ be the vertices with
$L_\phi(v)=\varnothing$.  For $v\notin S\cup U$, define
\[
 \delta_\phi(v)=
 \min_{i\in L_\phi(v)}
 \bigl|\{u\in N(v)\setminus(S\cup U):i\in L_\phi(u)\}\bigr|,
\]
and call $v$ restricting when $\delta_\phi(v)\geq\varepsilon n$; write $W$
for the restricting set.  Claim 17.2.2 constructs a total colouring after
deleting at most
\[
 (n-1)|S\cup U|+
 \frac12\sum_{v\notin S\cup U}\delta_\phi(v)
\]
monochromatic edges.  Since $G$ is $\varepsilon$-far, separating the terms
from $W$ and its complement gives Claim 17.2.3:
\[
 |S\cup U\cup W|\geq \frac{\varepsilon n}{2}.
\]

The proof now grows a ternary restriction tree.  A live leaf records one
partial colouring.  Hitting $U$ makes all three children terminal; hitting
$W$ branches over the at most three feasible colours and decreases the total
list mass by at least $\varepsilon n$.  Because that mass begins at $3n$,
every branch has depth at most $3/\varepsilon$ (Claim 17.2.4).  During the
$s=40/\varepsilon^2$ samples, $|S|\leq s\leq\varepsilon n/10$, so every fixed
live node advances in each round with conditional probability at least
$0.4\varepsilon$.  Thus its failure probability is bounded by
\[
 \Pr\!\left[\operatorname{Bin}
   \left(\frac{40}{\varepsilon^2},0.4\varepsilon\right)
   \leq\frac3\varepsilon\right]
 \leq \exp\!\left(-\frac{169}{32\varepsilon}\right).
\]
There are fewer than $3^{3/\varepsilon+1}$ possible tree nodes, so the union
bound is below $0.1$.  With probability at least $0.9$ every leaf is terminal;
Claim 17.2.5 then says that the sampled induced graph has no proper
3-colouring.  This is the exact role of Claims 17.2.2--17.2.6: edit distance
forces progress, list mass bounds depth, and the binomial tail completes all
branches simultaneously.

\subsection{Triangle-freeness and the removal lemma}

The triangle-removal lemma says that for every $\varepsilon>0$ there is
$\delta(\varepsilon)>0$ such that any $n$-vertex graph that is
$\varepsilon$-far from triangle-free contains at least
$\delta(\varepsilon)n^3$ triangles.  Its proof uses Szemerédi regularity:
delete edges incident with the exceptional class, inside parts, in irregular
pairs, and in low-density regular pairs.  If fewer than $\varepsilon n^2$
edges were deleted, one triangle remains among three dense regular pairs; the
counting argument then produces $\delta n^3$ triangles in the original graph.

The tester samples uniform triples and rejects if one spans a triangle.
It never rejects a triangle-free graph.  If $G$ is far, a random ordered triple
is a triangle with probability $\Omega(\delta)$, so
$O(1/\delta(\varepsilon))$ independent trials give constant rejection
probability.  Equivalently, sample a constant-size vertex set and query all
pairs.  The dependence on $\varepsilon$ is enormous because the regularity
lemma is quantitative, but it is independent of $n$.

\subsection{Hereditary properties}

A graph property is hereditary if it is closed under deleting vertices;
equivalently, it is characterized by a possibly infinite collection of forbidden
induced subgraphs.  Perfectness, chordality, and being an interval graph are
hereditary; a hereditary property need not be monotone under edge deletion.

The induced-removal theorem implies that every hereditary graph property is
testable with one-sided error in the dense model.  For every $\varepsilon>0$
there is $q=q_{\mathcal P}(\varepsilon)$ such that if $G$ is
$\varepsilon$-far from $\mathcal P$, then a uniform induced $q$-vertex subgraph
violates $\mathcal P$ with constant probability.  The tester samples $q$
vertices, queries all $\binom q2$ pairs, and checks the induced graph.
The distinction between ordinary and induced copies is essential:
edge-monotone properties use ordinary removal; hereditary properties require
induced removal.
\endgroup

\section*{Oral-exam decision map and traps}

\begin{center}\small
\begin{tabularx}{\textwidth}{@{}lX@{}}
\toprule
\textbf{If asked} & \textbf{Start with}\\
\midrule
Why two events correlate & lattice order, monotonicity, log-supermodular measure,
then FKG\\
Why a graph statistic concentrates & exposure/gradation, Doob martingale,
bounded difference, Azuma\\
Why Azuma is weak & redesign exposure or use a small certificate/Talagrand\\
Why a subgraph count concentrates & express it as a positive polynomial and
compute Kim--Vu derivative expectations\\
Where $1-H(p)$ comes from & exponential Hamming-ball volume and random-code
union bound\\
How randomness proves an adversarial lower bound & fix the opponent's strategy,
randomize the response, and use expectation/potential\\
How entropy proves counting & choose a uniform object, apply Shearer, bound
projection entropies by log support sizes\\
Why constant graph queries suffice & far implies many constant witnesses via a
removal/small-witness theorem\\
\bottomrule
\end{tabularx}
\end{center}

\begin{enumerate}
\item FKG requires both the correct partial order and a log-supermodular
measure; product measure is a special case, not the definition.
\item A martingale is not automatically concentrated: bounded differences or
another tail hypothesis must be proved.
\item Talagrand is centered naturally around a median/level set; translating to
the mean needs an additional argument.
\item The sphere bound is necessary for liar games/codes; existence at the
bound needs a separate construction.
\item Shannon's channel-coding radius is near $pn$ typical errors, not a
worst-case distance guarantee against every error pattern.
\item Property testing distinguishes graphs in $\mathcal P$ from graphs
$\varepsilon$-far from $\mathcal P$; behavior on graphs that are close but not
in $\mathcal P$ is unrestricted.
\item One-sided rejection must exhibit a real witness; for hereditary
properties the witness is an induced forbidden graph.
\end{enumerate}

\section*{Primary reference}
N.~Alon and J.~H.~Spencer, \emph{The Probabilistic Method}, 4th edition,
Wiley, 2016: Chapters 6, 7, 15, and 17 only.

\end{document}